\section{Vector flows and entropy on a random domain}
\label{sec:c-random-domains}

The next proof packages the conditional laws of all remaining cards
in one vector array. Averaging the next label's squared error cancels
the number of available positions in the entropy bound. As in
Proposition~\ref{lifts:cycle-entropy}, this gives a mean-information
estimate by an alternative to the uniform fixed-list argument
\eqref{noise:uniform-entropy}. We record it as a restriction to the
complement of a random omitted set. The application uses a different
Fourier transfer: its well-controlled domains may depend on the
sampled permutation.

As before $V=\mathbb F_2^d$, $n=2^d$, and layer $s$ switches along
$i_s=1+(s\bmod d)$.  Fix $0<\alpha\le1$ and an integer
$\alpha n\le m\le n$.  Choose a uniform $m$-subset $J$ of input
labels independently of the shuffle.  Reveal the paths of all labels
outside $J$.  At time $s$ let $O_s$ be the available positions, not
occupied by those revealed labels.  For $x\in V$, let $Q_x(s)\in\mathbb R^J$
have coordinates
\[
 Q_x(s)_b=\mathbb P\{\Pi_s(b)=x\mid J,
                  \hbox{paths of labels outside }J\hbox{ through }s\},
 \qquad b\in J,
\]
and define
\[
 X_x(s)=Q_x(s)-\frac{{\bf1}_{O_s}(x)}m\,\mathbf1_J.
\]
For each $b\in J$, the coordinate $X_x(s)_b$ is the centered tag law
$f_s(x)$ of Section~\ref{sec:first-route}, with the labels outside $J$
as blockers and $O_s$ as its available set.  Thus each coordinate sums
to zero over $x$.  Applying Lemma~\ref{lem:marginal-averaging} to each
tag gives the vector update: two available positions average their
vectors; one available position transports its vector to the unique
available output; and a pair with no available position carries zero.
The lemma also identifies these vectors when all exposed paths through
a terminal time are prescribed, and shows that the vector at an earlier
time depends only on their prefixes.  The total squared energy starts
at $m-1$ and is nonincreasing, so it is always at most $n$.

Let $\mathcal F_s$ contain $J$ and every actual switch coin used before
layer $s$.  This is the filtration used for expectations in the next
argument; it is distinct from the smaller sigma-field that defines the
conditional probabilities $Q_x(s)$.  Write
\[
 A_s=\mathbb E\sum_x\|X_x(s)\|_2^2,
 \qquad
 S_s=\mathbb E\sum_x\|X_x(s)\|_2^2
                  {\bf1}_{\{x+e_{i_s}\notin O_s\}}.
\]
The latter is the expected energy on edges with precisely one available
position.

\begin{proposition}[Vector noise and random-domain balance]
\label{prop:c-vector-domain-energy}
For $t\ge d$,
\begin{equation}
 A_t=\sum_{\ell=1}^d2^{-\ell}S_{t-\ell}.
 \label{eq:c-vector-noise-identity}
\end{equation}
There are constants $c_\alpha,C_\alpha>0$, depending only on $\alpha$,
such that, with $\gamma=\alpha/2$ and
$b_n=C_\alpha n e^{-c_\alpha n}$,
\begin{equation}
 S_s\le(1-\gamma)A_s+b_n\qquad(s\ge d-1).
 \label{eq:c-vector-singles-gain}
\end{equation}
In particular, for $\theta=1-\gamma/4$,
\begin{equation}
 A_t\le n\theta^{t-2d}+b_n/\gamma\qquad(t\ge0).
 \label{eq:c-vector-decay}
\end{equation}
Consequently there is a fixed integer $C>2$, depending only on $\alpha$,
for which $A_{Cd}=O(n^{-4})$, uniformly in $m\in[\alpha n,n]$.
\end{proposition}

\begin{proof}
For each $b\in J$, apply Lemma~\ref{noise:memory} to the tag $b$,
with blockers $V\setminus J$. Sum its identity over $b$ and then
average over $J$. The sum of the tag energies is exactly $A_t$, and
the sum of their weighted blocker occupancies is $S_s$, proving
\eqref{eq:c-vector-noise-identity}. Correlations between the tags
cause no cross terms: the squared array norm is the sum of their
individual squared norms. The new estimate needed here is therefore
the balance bound for this total energy.

To estimate the recent noise, fix $s\ge d-1$ and condition on
$\mathcal F_a$ at $a=s-d+1$.  The intervening $d-1$ layers act
independently in the two halves determined by coordinate $i_s$.
For $x$ in one half, $X_x(s)$ depends only on that half's subsequent
coin array and its initial vector array.  The event
$x+e_{i_s}\in O_s$ depends only on the other half's array and its
initial available set.  These quantities are therefore conditionally
independent.  Moreover, the latter conditional probability is exactly
the initial density of $O_a$ in the other half: averaging expected
occupancies through all its $d-1$ directions makes them constant there.

Unconditionally $O_a$ is a uniform $m$-set.  Indeed it is the image of
the independent uniform set $J$ under the actual permutation through
time $a$.  Sampling concentration gives that both half-densities are
at least $\alpha/2$, except on an $\mathcal F_a$-measurable event of
probability $O(e^{-c_\alpha n})$.  On its complement, conditional
independence gives available-partner energy at least $\gamma$ times
the same-time total energy.  On the exceptional event the energy is
at most $n$.  Taking expectations proves
\eqref{eq:c-vector-singles-gain}.  The factor $n$ in $b_n$ is needed
because this is the sum of the energies of all $m$ cards.

It remains to solve the recurrence.  Through time $2d$,
\eqref{eq:c-vector-decay} follows from $A_t\le n$.  At later times
substitute the induction hypothesis and
\eqref{eq:c-vector-singles-gain} into
\eqref{eq:c-vector-noise-identity}.  The coefficient of the proposed
geometric term is at most
\[
 (1-\gamma)\sum_{\ell\ge1}(2\theta)^{-\ell}
 =\frac{1-\gamma}{2\theta-1}<1,
\]
and the constant contribution is at most
$(1-\gamma)b_n/\gamma+b_n=b_n/\gamma$.
This proves \eqref{eq:c-vector-decay}.  Finally choose a fixed integer
$C>2$ large enough that $1+(C-2)\log_2\theta\le-4$.
Since $n=2^d$, its first term at $Cd$ is at most $n^{-4}$, and its
second term is exponentially small.
\end{proof}

\begin{theorem}[Mean relative entropy over omitted domains]
\label{thm:c-random-domain-entropy}
Let $d\ge8$, $q=256$, and let $R$ be a uniform $(n/q)$-subset of
input labels, independent of the shuffle.  For a permutation $g$, write
$g_{-R}$ for its restriction to $V\setminus R$, retaining the identities
of all those labels.  There is an absolute positive integer $C$ such
that, for the law $\mu$ after $Cd$ physical shuffles and the uniform
restriction law $u_{-R}$,
\begin{equation}
 \mathbb E_R D(\mu_{-R}\Vert u_{-R})=O(n^{-3})=o(1).
 \label{eq:c-random-domain-entropy}
\end{equation}
Relative entropy uses natural logarithms.
\end{theorem}

\begin{proof}
Apply Proposition~\ref{prop:c-vector-domain-energy} with
$\alpha=1/256$. Generate the omitted set and the retained order
together: choose one uniform order $(c_1,\ldots,c_n)$, independently
of the shuffle, and let $R$ consist of its last $n/256$ labels.
For every fixed order, the first $n-n/256$ endpoints encode the
restriction outside $R$, so their relative entropy is
$D(\mu_{-R}\Vert u_{-R})$.

At a rank with $j$ earlier labels, the entropy chain's reference law
is uniform on the $m=n-j$ positions not occupied by their endpoints.
Conditioning further on those labels' whole paths increases the
expected contribution, by convexity, because the reference set is
already determined by the endpoints. Put
$J=V\setminus\{c_1,\ldots,c_j\}$ and $b=c_{j+1}$.
For fixed preceding labels and a fixed next label $b$, this
path-conditioned endpoint law is $Q(T)_b$. It does not depend on
the order of the later labels or on which of them form $R$: these
names contain no information about the switches.

We may therefore average the unrevealed names before taking the
next chain contribution. Conditional only on the preceding label
choices and their paths, $b$ is uniform in $J$. This averaging step
does not fix $R$; if it did, $b$ would be uniform only in $J\setminus R$.
For each possible $b$, the elementary inequality
$\log z\le z-1$ yields the chi-square bound
\[
 D\bigl(Q(T)_b\Vert U_{O_T}\bigr)
 \le m\sum_{x\in O_T}|X_x(T)_b|^2.
\]
Averaging this display over $b\in J$ gives
$\sum_x\|X_x(T)\|_2^2$. Before the preceding labels are fixed,
$J$ is a uniform $m$-set independent of the shuffle, so its remaining
expectation is $A_T$. Every stage has $m\ge n/256$, and hence its
expected contribution is $O(n^{-4})$ uniformly. There are at most
$n$ stages. Averaging the entropy chain over the full order proves
\eqref{eq:c-random-domain-entropy}.
\end{proof}

The transfer below needs only the mean over $R$, although the earlier
bound \eqref{noise:uniform-entropy} also gives small entropy for every
fixed omitted domain of this size after sufficiently many sweeps.
Its use of the mean permits the domains having bounded density to
depend on the sampled permutation.

\begin{corollary}[Random-domain application]
\label{cor:c-random-domain-mixing}
With the fixed integer $C$ from
Theorem~\ref{thm:c-random-domain-entropy}, the full permutation law after
$32Cd$ physical shuffles tends to uniform in total variation from
every starting deck.
\end{corollary}

\begin{proof}
The random-domain transfer
\cite[Theorem~\ref*{l-l:domain-transfer}]{partialFourier}
takes $q=256$ and the mean entropy bound
\eqref{eq:c-random-domain-entropy}.  Its truncation retains permutations
$g$ for which at least a fraction $7/8$ of the domains satisfy
$\mu_{-R}(g_{-R})/u_{-R}(g_{-R})\le2$.  It gives a probability $\nu$
with $\|\nu-\mu\|_{\rm TV}=o(1)$ and
\[
 \|\widehat\nu(\lambda)\|_{\rm op}
 \le C_{\rm dom}D_\lambda^{-1/8}\qquad(D_\lambda>1),
\]
where $C_{\rm dom}$ is absolute.  Thus the mean-domain hypothesis, the fraction
$7/8$, and the density cutoff two are precisely the ones supplied to
that theorem; simultaneous control of every $R$ is unnecessary.

Apply the probability completion of
Section~\ref{sec:probability-completion}, with operator exponent
$1/8$ and thirty-two factors. The original block's zero sign mean
and $\|\nu-\mu\|_{\rm TV}=o(1)$ supply the scalar condition.
The nonsign Plancherel sum is bounded by
\[
 C_{\rm dom}^{64}\sum_{D_\lambda>1}D_\lambda^{2-64/8}
 =C_{\rm dom}^{64}\sum_{D_\lambda>1}D_\lambda^{-6}=o(1),
\]
by the inverse-square degree estimate
\cite[Lemma~\ref*{l-l:degree-inverse-square}]{partialFourier}.
The completion therefore gives convergence of the original law at
$32Cd$ physical shuffles, uniformly over deterministic starting decks.
\end{proof}
